% Part: normal-modal-logic
% Chapter: tableaux
% Section: rules-for-K
%READERNOTE{SOL6-A175: नव्या पूर्वचिन्हाच्या अटीचे प्रतिउदाहरण दाखवणाऱ्या केवळ Box-पर्यायात असत्य Box सूत्रावर लागू होणाऱ्या नियमाचे लेबल असत्य Box असे दुरुस्त केले. उदाहरणातील जाणीवपूर्वक शिथिल केलेली नवेपणाची अट बदललेली नाही.}

\documentclass[../../../include/open-logic-section]{subfiles}

\begin{document}

\olfileid{nml}{tab}{rul}

\olsection{\Ax{K} साठीचे नियम}

नेहमीच्या विधानीय कारकांचे नियम नेहमीच्या विधानीय
चिन्हांकित !!{tableau}s मधील नियमांसारखेच आहेत;
फक्त त्यांना पूर्वचिन्हे जोडली जातात. प्रत्येक वेळी,
$\sFmla{S}{!A}[\sigma]$ या चिन्हांकित !!{formula}
ला नियम लावल्याने मिळणाऱ्या नव्या !!{formula}s
वरही~$\sigma$ हेच पूर्वचिन्ह असते. हे सहज लक्षात
येते: उदाहरणार्थ, $!A \land !B$ हे~$\sigma$ ने
नाव दिलेल्या जगात सत्य असेल, तर $!A$ आणि $!B$
ही दोन्ही~$\sigma$ येथे सत्य असतात (इतर कोणत्याही
जगात असतीलच असे नाही). विधानीय नियम
\olref{tab:prop-rules} मध्ये दिले आहेत.

\begin{table}
  \[\def\arraystretch{3}\begin{array}{|c|c|}
    \hline
    \AxiomC{\sFmla{\True}{\lnot !A}[\sigma]}
    \RightLabel{\TRule{\True}{\lnot}}
    \UnaryInfC{\sFmla{\False}{!A}[\sigma]}
    \DisplayProof
    &
    \AxiomC{\sFmla{\False}{\lnot !A}[\sigma]}
    \RightLabel{\TRule{\False}{\lnot}}
    \UnaryInfC{\sFmla{\True}{!A}[\sigma]}
    \DisplayProof
    \\[1ex]
    \hline
    \AxiomC{\sFmla{\True}{!A \land !B}[\sigma]}
    \RightLabel{\TRule{\True}{\land}}
    \UnaryInfC{\sFmla{\True}{!A}[\sigma]}
    \noLine
    \UnaryInfC{\sFmla{\True}{!B}[\sigma]}
    \DisplayProof
    &
    \AxiomC{\sFmla{\False}{!A \land !B}[\sigma]}
    \RightLabel{\TRule{\False}{\land}}
    \UnaryInfC{$\sFmla{\False}{!A}[\sigma] \quad \mid \quad
      \sFmla{\False}{!B}[\sigma]$}
    \DisplayProof
    \\[2ex]
    \hline
    \AxiomC{\sFmla{\True}{!A \lor !B}[\sigma]}
    \RightLabel{\TRule{\True}{\lor}}
    \UnaryInfC{$\sFmla{\True}{!A}[\sigma] \quad \mid \quad
      \sFmla{\True}{!B}[\sigma]$}
    \DisplayProof
    &
    \AxiomC{\sFmla{\False}{!A \lor !B}[\sigma]}
    \RightLabel{\TRule{\False}{\lor}}
    \UnaryInfC{\sFmla{\False}{!A}[\sigma]}
    \noLine
    \UnaryInfC{\sFmla{\False}{!B}[\sigma]}
    \DisplayProof
    \\[2ex]
    \hline
    \AxiomC{\sFmla{\True}{!A \lif !B}[\sigma]}
    \RightLabel{\TRule{\True}{\lif}}
    \UnaryInfC{$\sFmla{\False}{!A}[\sigma] \quad \mid
      \quad \sFmla{\True}{!B}[\sigma]$}
    \DisplayProof
    &
    \AxiomC{\sFmla{\False}{!A \lif !B}[\sigma]}
    \RightLabel{\TRule{\False}{\lif}}
    \UnaryInfC{\sFmla{\True}{!A}[\sigma]}
    \noLine
    \UnaryInfC{\sFmla{\False}{!B}[\sigma]}
    \DisplayProof
    \\[2ex]
    \hline
  \end{array}\]
  \caption{विधानीय कारकांसाठी पूर्वचिन्हयुक्त
    !!{tableau} नियम}
  \ollabel{tab:prop-rules}
\end{table}

शाखा बंद होण्याची अट नेहमीच्या !!{tableau}s
प्रमाणेच आहे; पण येथे !!{formula}s बरोबर त्यांची
पूर्वचिन्हेही जुळली पाहिजेत. म्हणजे एखाद्या शाखेत
पुढील दोन्ही असतील, तर ती बंद असते:
\[
\sFmla{\True}{!A}[\sigma] \quad\text{आणि}\quad \sFmla{\False}{!A}[\sigma]
\]
येथे $\sigma$ हे एखादे पूर्वचिन्ह आणि $!A$ हे
!!{formula} आहे.

गृहीतके मांडण्याचे नियमही नेहमीच्या !!{tableau}s
प्रमाणेच आहेत; मात्र गृहीतकांसाठी आपण नेहमी
$1$ हे पूर्वचिन्ह वापरतो. (सर्व गृहीतकांसाठी एकच
पूर्वचिन्ह वापरले, तर नेमके कोणते पूर्वचिन्ह
वापरले याने फरक पडत नाही.) म्हणून, उदाहरणार्थ,
\[
!B_1, \dots, !B_n \Proves !A
\]
असे तेव्हा आणि तेव्हाच म्हणतो जेव्हा पुढील
गृहीतकांसाठी बंद टॅब्लो असतो:
\[
\sFmla{\True}{!B_1}[1], \dots, \sFmla{\True}{!B_n}[1],
\sFmla{\False}{!A}[1].
\]

मोडल कारक\iftag{prvBox}{\iftag{prvDiamond}{~$\Box$
    आणि}{~$\Box$}}{}\iftag{prvDiamond}{~$\Diamond$}{}
यांसाठी,~$\sigma$ हे पूर्वचिन्ह असलेल्या !!a{formula}
ला नियम लावल्यावर निष्कर्षाचे पूर्वचिन्ह
$\sigma.n$ असते. पण कोणता $n$ चालेल हे
चिन्ह~$\True$ आहे की~$\False$ यावर अवलंबून असते.

\iftag{prvBox}{$\TRule{\True}{\Box}$ हा नियम
  $\sFmla{\True}{\Box !A}[\sigma]$ असलेल्या शाखेत
  $\sFmla{\True}{!A}[\sigma.n]$ जोडतो.\iftag{prvDiamond}{
    त्याचप्रमाणे, }{}}{}\iftag{prvDiamond}{$\TRule{\False}{\Diamond}$
  हा नियम $\sFmla{\False}{\Diamond !A}[\sigma]$
  असलेल्या शाखेत $\sFmla{\False}{!A}[\sigma.n]$
  जोडतो.}{}
\iftag{notprvBox,notprvDiamond}{हा नियम}{हे नियम}
लावताना $\sigma.n$ हे पूर्वचिन्ह संबंधित शाखेत
\emph{आधीपासूनच} आलेले असले पाहिजे. अशा
पूर्वचिन्हाला त्या शाखेवर ‘वापरलेले’ म्हणू.

\iftag{prvBox}{$\TRule{\False}{\Box}$ हा नियम
  $\sFmla{\False}{\Box !A}[\sigma]$ असलेल्या शाखेत
  $\sFmla{\False}{!A}[\sigma.n]$ जोडतो.\iftag{prvDiamond}{
    त्याचप्रमाणे, }{}}{}\iftag{prvDiamond}{$\TRule{\True}{\Diamond}$
  हा नियम $\sFmla{\True}{\Diamond !A}[\sigma]$
  असलेल्या शाखेत $\sFmla{\True}{!A}[\sigma.n]$
  जोडतो.}{}
परंतु \iftag{notprvBox,notprvDiamond}{हा नियम}{हे नियम}
लावताना $\sigma.n$ हे पूर्वचिन्ह संबंधित शाखेत
\emph{आधीपासून आलेले नसले} पाहिजे. अशा
पूर्वचिन्हांना त्या शाखेसाठी ‘नवी’ म्हणतो.

हे नियम \olref{tab:rules-K} मध्ये दिले आहेत.

\begin{table}
  \begin{center}
    \def\arraystretch{3}\def\fCenter{}
    \begin{tabular}{|c|c|}
    \hline
    \iftag{prvBox}{
      \AxiomC{\sFmla{\True}{\Box !A}[\sigma]}
      \RightLabel{\TRule{\True}{\Box}}
      \UnaryInfC{\sFmla{\True}{!A}[\sigma.n]}
      \DisplayProof
      &
      \AxiomC{\sFmla{\False}{\Box !A}[\sigma]}
      \RightLabel{\TRule{\False}{\Box}}
      \UnaryInfC{\sFmla{\False}{!A}[\sigma.n]}
      \DisplayProof\\
      $\sigma.n$ वापरलेले आहे & $\sigma.n$ नवे आहे
      \\[1ex]
      \hline}{}
    \iftag{prvDiamond}{
      \AxiomC{\sFmla{\True}{\Diamond !A}[\sigma]}
      \RightLabel{\TRule{\True}{\Diamond}}
      \UnaryInfC{\sFmla{\True}{!A}[\sigma.n]}
      \DisplayProof
      &
      \AxiomC{\sFmla{\False}{\Diamond !A}[\sigma]}
      \RightLabel{\TRule{\False}{\Diamond}}
      \UnaryInfC{\sFmla{\False}{!A}[\sigma.n]}
      \DisplayProof\\
      $\sigma.n$ नवे आहे & $\sigma.n$ वापरलेले आहे
      \\[1ex]
      \hline}{}
    \end{tabular}
  \end{center}
  \caption{\Ax{K} साठीचे मोडल नियम.}
  \ollabel{tab:rules-K}
\end{table}

\iftag{prvBox}{\TRule{\True}{\Box}}{\TRule{\False}{\Diamond}}
या नियमासाठी पूर्वचिन्ह वापरलेले असणे आवश्यक आहे.
ही अट नसेल, तर पुढील !!{tableau} बंद असल्याचे
आपण चुकीने मानू:
\iftag{notprvBox,notprvDiamond}{%
  \iftag{prvBox}{%
  \begin{oltableau}
    [\pFmla{\True}{\Box \formula{A}}{1}, just = \TAss
      [\pFmla{\False}{\lnot\Box\lnot \formula{A}}{1}, just = \TAss
        [\pFmla{\True}{\formula{A}}{1.1}, just = {\TRule{\True}{\Box}[1]}
          [\pFmla{\True}{\Box\lnot\formula{A}}{1},
            just ={\TRule{\False}{\lnot}[2]}
            [\pFmla{\True}{\lnot\formula{A}}{1.1},
              just = {\TRule{\True}{\Box}[4]}
              [\pFmla{\False}{\formula{A}}{1.1},
                  just ={\TRule{\True}{\lnot}[5]}, close]
            ]
          ]
        ]
      ]
    ]
  \end{oltableau}
  }{
  \begin{oltableau}
    [\pFmla{\True}{\lnot\Diamond\lnot \formula{A}}{1}, just = \TAss
      [\pFmla{\False}{\Diamond\formula{A}}{1}, just = \TAss
        [\pFmla{\False}{\formula{A}}{1.1}, just = {\TRule{\False}{\Diamond}[2]}
          [\pFmla{\False}{\Diamond\lnot\formula{A}}{1},
            just ={\TRule{\True}{\lnot}[1]}
            [\pFmla{\False}{\lnot\formula{A}}{1.1},
              just = {\TRule{\False}{\Diamond}[4]}
              [\pFmla{\True}{\formula{A}}{1.1},
                  just ={\TRule{\True}{\lnot}[5]}, close]
            ]
          ]
        ]
      ]
    ]
  \end{oltableau}
}}{
  \begin{oltableau}
    [\pFmla{\True}{\Box \formula{A}}{1}, just = \TAss
      [\pFmla{\False}{\Diamond \formula{A}}{1}, just = \TAss
        [\pFmla{\True}{\formula{A}}{1.1}, just = {\TRule{\True}{\Box}[1]}
          [\pFmla{\False}{\formula{A}}{1.1},
            just = {\TRule{\False}{\Diamond}[2]}, close]
        ]
      ]
    ]
\end{oltableau}}

पण $\Box \formula{A} \Entails/ \Diamond \formula{A}$;
म्हणून असे झाल्यास आपली प्रमाणपद्धती अप्रमाण असेल.
त्याचप्रमाणे, $\Diamond \formula{A} \Entails/
\Box \formula{A}$; परंतु
\iftag{prvBox}{\TRule{\False}{\Box}}{\TRule{\True}{\Diamond}}
या नियमासाठी पूर्वचिन्ह नवे असण्याची अट नसेल,
तर पुढील टॅब्लो बंद होईल: \iftag{notprvBox,notprvDiamond}{%
  \iftag{prvBox}{%
  \begin{oltableau}
    [\pFmla{\True}{\lnot\Box\lnot \formula{A}}{1}, just = \TAss
      [\pFmla{\False}{\Box \formula{A}}{1}, just = \TAss
        [\pFmla{\False}{\formula{A}}{1.1}, just = {\TRule{\False}{\Box}[2]}
          [\pFmla{\False}{\Box\lnot\formula{A}}{1},
            just ={\TRule{\True}{\lnot}[1]}
            [\pFmla{\False}{\lnot\formula{A}}{1.1},
              just = {\TRule{\False}{\Box}[4]}
              [\pFmla{\True}{\formula{A}}{1.1},
                  just ={\TRule{\False}{\lnot}[5]}, close]
            ]
          ]
        ]
      ]
    ]
  \end{oltableau}
  }{
  \begin{oltableau}
    [\pFmla{\True}{\Diamond \formula{A}}{1}, just = \TAss
      [\pFmla{\False}{\lnot\Diamond\lnot\formula{A}}{1}, just = \TAss
        [\pFmla{\True}{\formula{A}}{1.1}, just = {\TRule{\True}{\Diamond}[1]}
          [\pFmla{\True}{\Diamond\lnot\formula{A}}{1},
            just ={\TRule{\False}{\lnot}[2]}
            [\pFmla{\True}{\lnot\formula{A}}{1.1},
              just = {\TRule{\True}{\Diamond}[4]}
              [\pFmla{\False}{\formula{A}}{1.1},
                  just ={\TRule{\True}{\lnot}[5]}, close]
            ]
          ]
        ]
      ]
    ]
  \end{oltableau}
}}{
  \begin{oltableau}
    [\pFmla{\True}{\Diamond \formula{A}}{1}, just = \TAss,name=one
      [\pFmla{\False}{\Box \formula{A}}{1}, just = \TAss, name=two
        [\pFmla{\True}{\formula{A}}{1.1}, just = {\TRule{\True}{\Diamond}[1]}
          [\pFmla{\False}{\formula{A}}{1.1},
            just = {\TRule{\False}{\Box}[2]}, close]
        ]
      ]
    ]
  \end{oltableau}}

\end{document}
